%% ============================================================
\section{A Scaling Law for When Adaptation Pays}
\label{sec:law}

The RTT sweep of Section~\ref{sec:results} left a result we could not explain.
Adaptation recovered a large, highly significant share of bulk-flow round-trip
time on a $5$\,ms path and nothing at $20$, $80$ or $200$\,ms. Read as a table
that is four measurements and a shrug. The obvious reading, that adaptation
helps on short paths, is not a mechanism and predicts nothing.

\subsection{The hypothesis}

CoDel's \texttt{target} is a delay budget, fixed at $5$\,ms by default, while
the path RTT is whatever the path is. The quantity that should govern whether
that budget is badly chosen is therefore neither the target nor the RTT but
their ratio,
\begin{equation}
r = \frac{\texttt{target}}{\text{path RTT}} .
\end{equation}
When $r$ is small the standing-queue budget is a minor addition to the path and
there is nothing for a controller to win back. When $r$ approaches unity the
budget is comparable to the path itself, queueing dominates the delay a flow
sees, and shrinking it is worth a great deal. On this reading RTT never appears
on its own; it appears only through $r$. That is a stronger claim than the
table, and unlike the table it can be wrong.

\subsection{Pre-registration}

Because a saturating curve can be fitted to almost any monotone data, the
hypothesis, the three predictions and their numerical decision thresholds were
written down and committed to the repository \emph{before} the measurement
campaign ran. The commit adding \texttt{docs/SCALING\_LAW.md} and the analysis
script, with its thresholds already in place, precedes the commit adding the
results directory. The predictions were:

\begin{description}
\item[P1] A crossover exists between $r = 0.25$ and $r = 1.0$.
\item[P2] \textbf{Ratio invariance.} Configurations sharing an $r$ but
      differing in RTT show the same benefit. Threshold: benefits within $10$
      percentage points. This is the prediction that can refute the hypothesis
      outright, and a refutation was to be reported as such.
\item[P3] Leave-one-out prediction across interior ratios, with mean absolute
      error at or below $5$ percentage points.
\end{description}

One amendment was made during the campaign and is recorded rather than
substituted silently. The first attempt at P2 varied the target while leaving
\texttt{interval} at $100$\,ms, which moved \texttt{target}/\texttt{interval}
from $0.05$ to $0.80$ across the cells and so varied a second dimensionless
group alongside the one under test. The interval is now scaled with the target,
and the claim is correspondingly narrower: \emph{at fixed}
\texttt{target}/\texttt{interval}, the benefit is governed by $r$.

\subsection{Results}

Table~\ref{tab:law} reports every cell. Figure~\ref{fig:law} plots benefit
against $r$ and, in the second panel, the same data against RTT alone.

\begin{table}[t]
\centering
\caption{Adaptation benefit against the ratio $r = \texttt{target}/\text{RTT}$.
Reduction in mean bulk-flow RTT, static against adapted \texttt{fq\_codel}.
Starred rows are significant at $p < 0.05$. \emph{adj} is the mean number of
parameter adjustments the controller made; a cell where it is zero is not
measuring adaptation and is excluded from the fit.}
\label{tab:law}
\small
% GENERATED by src/analyse_law.py -- do not edit by hand.
\begin{tabular}{rrrrrrrl}
\hline
target & RTT & $r$ & static & adapted & \multicolumn{2}{c}{reduction} & $p$ \\
 (ms) & (ms) & & (ms) & (ms) & (\%) & (ms) & \\
\hline
5 & 2 & 2.500 & 26.10 & 21.10 & 19.2$^{*}$ & 5.00 & $<$0.001 \\
5 & 5 & 1.000 & 27.10 & 22.10 & 18.5$^{*}$ & 5.00 & $<$0.001 \\
20 & 20 & 1.000 & 40.10 & 32.60 & 18.7$^{*}$ & 7.50 & $<$0.001 \\
5 & 20 & 0.250 & 36.10 & 35.60 & 1.4$^{*}$ & 0.50 & $<$0.001 \\
20 & 80 & 0.250 & 92.10 & 72.10 & 21.7$^{*}$ & 20.00 & $<$0.001 \\
5 & 80 & 0.062 & 89.10 & 89.60 & -0.6$^{*}$ & -0.50 & $<$0.001 \\
5 & 200 & 0.025 & 208.10 & 207.60 & 0.2$^{*}$ & 0.50 & $<$0.001 \\
\hline
\end{tabular}

\end{table}

\begin{figure}[t]
\centering
\includegraphics[width=\columnwidth]{fig16_scaling_law.png}
\caption{Benefit against $r$ (left) and against RTT alone (right). The
collapse onto a single curve in the left panel, and its absence in the right,
is the result.}
\label{fig:law}
\end{figure}

\textbf{P1 holds.} The crossover falls where it was predicted to, between
$r = 0.25$ and $r = 1.0$, and the curve saturates rather than growing without
bound. Fitting $b(r) = B / (1 + (r_0/r)^k)$ gives
\begin{equation}
b(r) = \frac{\LawCeiling}{1 + (\LawHalfBenefit/r)^{\LawSteepness}}, \qquad
R^2 = \LawRsq .
\end{equation}
Half the available benefit arrives at $r = \LawHalfBenefit$. The significance
boundary coincides with the crossover: every cell at $r \ge \LawLowestSigRatio$
is significant and every cell below it is not.

A percentage reduction shrinks on its own when its denominator grows, so the
collapse at low $r$ could be an artefact of dividing by a larger RTT. It is
not. In absolute milliseconds the two regimes do not overlap: the gain is
$\LawGainHiMin$ to $\LawGainHiMax$\,ms above $r = 0.4$ and $\LawGainLoMin$ to
$\LawGainLoMax$\,ms below it.

\textbf{P2 holds, at $r = 1.0$, across a sixteenfold range of RTT.} Scaling the
target, the interval, the controller's loop period and the run duration
together by the same factor:

\begin{center}
\small
\begin{tabular}{rrrrr}
\hline
target & RTT & benefit & gain & $p$ \\
\hline
5\,ms & 5\,ms & 18.4\% & 5.10\,ms & $<$0.001 \\
20\,ms & 20\,ms & 17.6\% & 10.78\,ms & $<$0.001 \\
80\,ms & 80\,ms & 16.4\% & 29.73\,ms & 0.0003 \\
\hline
\end{tabular}
\end{center}

Spread is $2.0$ percentage points against a threshold of $10$, and all
three cells are significant. The \emph{fraction} of delay recovered is what
stays fixed while the absolute gain grows from $5$ to $30$\,ms, which is the
claim itself: the ratio governs the proportion recoverable, not the
milliseconds.

The sharpest discrimination was not anticipated in the pre-registration. Hold
RTT fixed at $20$\,ms and vary only $r$: the benefit is $1.4\%$ and not
significant at $r = 0.25$, and $17.6\%$ and significant at $r = 1.0$. Same
path, twelve times the benefit. The competing explanation, that short paths are
simply special, predicts nothing for both, since $20$\,ms is one of the RTTs at
which the original sweep measured nothing. The ratio predicts both the presence
and the absence; RTT predicts neither.

Invariance at $r = 0.25$ is reported as inconclusive rather than as
supporting evidence. One of its two cells made no parameter adjustment at all
across three seeds, because its drop rate sat almost exactly on the scaled
regime threshold and the classification oscillated. A cell whose adapted arm
never changes a parameter runs the static configuration twice, so its null is
true by construction and cannot distinguish an absent effect from an absent
intervention.

\textbf{P3 holds.} Leave-one-out across interior ratios predicts held-out
measurements to $1.4$ percentage points on average, worst case $3.5$, against a
threshold of $5$. Endpoints are excluded deliberately: holding out an extreme
ratio asks a saturating curve to extrapolate past the data that fixes its own
ceiling, which tests the fitting procedure rather than the hypothesis.

\subsection{What the benefit costs}

Adaptation is not free, and reporting the latency gain without its price
reports half the trade. Shrinking the delay target makes CoDel drop earlier, so
senders back off sooner. Across the cells measured, the drop rate rises by up
to $\LawDropRiseHigh\%$ and retransmissions by a comparable margin wherever the
benefit is large, while both fall slightly in the cells where it is absent.
Throughput is the quantity that does not move: the worst change across every
cell in the campaign is $\LawWorstThroughput\%$. The controller converts
retransmissions into latency at a roughly fixed exchange rate, and the ratio
governs the exchange as well as the gain.

\subsection{Consequence}

Published guidance for CoDel holds that \texttt{target} should be about $5$ to
$10$ percent of \texttt{interval}, with \texttt{interval} on the order of the
path RTT. A correctly configured deployment therefore sits at $r$ of roughly
$0.05$ to $0.10$. Half the available benefit arrives at $r = \LawHalfBenefit$,
five to ten times outside that envelope.

This says something sharper than that adaptation sometimes helps.
It says adaptation pays only on boxes already configured well outside CoDel's
own design envelope, and that on a correctly configured box there is close to
nothing for a controller to recover. This reading unifies the null result of
Section~\ref{sec:results} with the null results reported elsewhere in the
literature as a single effect rather than as separate disappointments, and it
gives an operator a number to compute before deploying a controller at all.
